% Tensor-diagram recipe, shared like preamble.tex. A page that draws one starts
%   \documentclass[border=4pt]{standalone}
%   % Shared by every page under figures/src/: each page is
%   \documentclass[border=4pt]{standalone}
%   % Shared by every page under figures/src/: each page is
%   \documentclass[border=4pt]{standalone}
%   % Shared by every page under figures/src/: each page is
%   \documentclass[border=4pt]{standalone}
%   \input{preamble}
% and build.sh (next to this file) puts this directory on TEXINPUTS so the
% \input resolves.
\usepackage{amsmath}
\usepackage{amssymb}
\usepackage{braket}
\usepackage{booktabs}
\usepackage{xcolor}
% Colours shared with the schematic pages (schematic.py).
\definecolor{nuc}{HTML}{C96A1B}
\definecolor{ele}{HTML}{2E6FBA}

% and build.sh (next to this file) puts this directory on TEXINPUTS so the
% \input resolves.
\usepackage{amsmath}
\usepackage{amssymb}
\usepackage{braket}
\usepackage{booktabs}
\usepackage{xcolor}
% Colours shared with the schematic pages (schematic.py).
\definecolor{nuc}{HTML}{C96A1B}
\definecolor{ele}{HTML}{2E6FBA}

% and build.sh (next to this file) puts this directory on TEXINPUTS so the
% \input resolves.
\usepackage{amsmath}
\usepackage{amssymb}
\usepackage{braket}
\usepackage{booktabs}
\usepackage{xcolor}
% Colours shared with the schematic pages (schematic.py).
\definecolor{nuc}{HTML}{C96A1B}
\definecolor{ele}{HTML}{2E6FBA}

%   % Tensor-diagram recipe, shared like preamble.tex. A page that draws one starts
%   \documentclass[border=4pt]{standalone}
%   % Shared by every page under figures/src/: each page is
%   \documentclass[border=4pt]{standalone}
%   % Shared by every page under figures/src/: each page is
%   \documentclass[border=4pt]{standalone}
%   \input{preamble}
% and build.sh (next to this file) puts this directory on TEXINPUTS so the
% \input resolves.
\usepackage{amsmath}
\usepackage{amssymb}
\usepackage{braket}
\usepackage{booktabs}
\usepackage{xcolor}
% Colours shared with the schematic pages (schematic.py).
\definecolor{nuc}{HTML}{C96A1B}
\definecolor{ele}{HTML}{2E6FBA}

% and build.sh (next to this file) puts this directory on TEXINPUTS so the
% \input resolves.
\usepackage{amsmath}
\usepackage{amssymb}
\usepackage{braket}
\usepackage{booktabs}
\usepackage{xcolor}
% Colours shared with the schematic pages (schematic.py).
\definecolor{nuc}{HTML}{C96A1B}
\definecolor{ele}{HTML}{2E6FBA}

%   % Tensor-diagram recipe, shared like preamble.tex. A page that draws one starts
%   \documentclass[border=4pt]{standalone}
%   % Shared by every page under figures/src/: each page is
%   \documentclass[border=4pt]{standalone}
%   \input{preamble}
% and build.sh (next to this file) puts this directory on TEXINPUTS so the
% \input resolves.
\usepackage{amsmath}
\usepackage{amssymb}
\usepackage{braket}
\usepackage{booktabs}
\usepackage{xcolor}
% Colours shared with the schematic pages (schematic.py).
\definecolor{nuc}{HTML}{C96A1B}
\definecolor{ele}{HTML}{2E6FBA}

%   % Tensor-diagram recipe, shared like preamble.tex. A page that draws one starts
%   \documentclass[border=4pt]{standalone}
%   \input{preamble}
%   \input{tensor}
%
% The convention: an expansion is a contraction, drawn as in MPS / SVD diagrams,
% except that a leg is one of two kinds.
%   cont  a wavy leg: a continuous argument, r or r_1 (a function of position)
%   disc  a plain line: a finite index, i or mu (ordinary linear algebra)
% Nodes likewise:
%   fn    a circle, in the electron colour: a function of position (psi, phi, chi)
%   coef  a square: a plain array of numbers (C, c)
%   coefwide a wide square box: an array with several legs on one side
%         (C_ij, H_ij,kl)
%   fntall, coeftall  the same two turned upright, for diagrams read left to
%         right (legs leave from the left and right sides)
%   frame a dashed box around a group of nodes that contracts to one array
%   fnwide a wide rounded box: a function of several positions (psi(r1, r2)),
%         its cont legs dropped in parallel from the bottom edge
% Joining two legs means summing over that index, so
%   psi(r1, r2) = sum_ij C_ij phi_i(r1) phi_j(r2)
% is  fn with two cont legs  =  fn -disc- coef -disc- fn,  each fn keeping its
% cont leg. Label a leg with node[leg] at its free end (or mid-way for a
% contracted one, e.g. "i").
%
% Minimal use:
%   \begin{tikzpicture}
%     \node[fn] (p) at (0,0) {$\varphi$};
%     \draw[cont] (p) -- ++(0,-1) node[leg, below] {$\mathbf r$};
%     \draw[disc] (p) -- ++(-1,0) node[leg, left] {$i$};
%   \end{tikzpicture}
\usepackage{tikz}
\usetikzlibrary{decorations.pathmorphing}
\tikzset{
  fn/.style   = {circle, draw=ele, fill=ele!12, line width=0.9pt,
                 minimum size=8mm, inner sep=0pt},
  fnwide/.style = {fn, rectangle, rounded corners=3.5mm, minimum width=17mm,
                   minimum height=8mm},
  coef/.style = {rectangle, draw=black!80, fill=black!6, line width=0.9pt,
                 minimum size=7mm, inner sep=0pt},
  coefwide/.style = {coef, minimum width=17mm, minimum height=7mm},
  fntall/.style = {fn, rectangle, rounded corners=3.5mm, minimum width=8mm,
                   minimum height=17mm},
  coeftall/.style = {coef, minimum width=7mm, minimum height=17mm},
  frame/.style = {draw=black!50, dashed, rounded corners=2mm, line width=0.7pt},
  cont/.style = {draw=ele, line width=0.9pt, decorate,
                 decoration={snake, amplitude=0.55mm, segment length=2.4mm,
                             pre length=0.4mm, post length=0.8mm}},
  disc/.style = {draw=black!80, line width=0.9pt},
  leg/.style  = {font=\small, inner sep=2pt},
}

%
% The convention: an expansion is a contraction, drawn as in MPS / SVD diagrams,
% except that a leg is one of two kinds.
%   cont  a wavy leg: a continuous argument, r or r_1 (a function of position)
%   disc  a plain line: a finite index, i or mu (ordinary linear algebra)
% Nodes likewise:
%   fn    a circle, in the electron colour: a function of position (psi, phi, chi)
%   coef  a square: a plain array of numbers (C, c)
%   coefwide a wide square box: an array with several legs on one side
%         (C_ij, H_ij,kl)
%   fntall, coeftall  the same two turned upright, for diagrams read left to
%         right (legs leave from the left and right sides)
%   frame a dashed box around a group of nodes that contracts to one array
%   fnwide a wide rounded box: a function of several positions (psi(r1, r2)),
%         its cont legs dropped in parallel from the bottom edge
% Joining two legs means summing over that index, so
%   psi(r1, r2) = sum_ij C_ij phi_i(r1) phi_j(r2)
% is  fn with two cont legs  =  fn -disc- coef -disc- fn,  each fn keeping its
% cont leg. Label a leg with node[leg] at its free end (or mid-way for a
% contracted one, e.g. "i").
%
% Minimal use:
%   \begin{tikzpicture}
%     \node[fn] (p) at (0,0) {$\varphi$};
%     \draw[cont] (p) -- ++(0,-1) node[leg, below] {$\mathbf r$};
%     \draw[disc] (p) -- ++(-1,0) node[leg, left] {$i$};
%   \end{tikzpicture}
\usepackage{tikz}
\usetikzlibrary{decorations.pathmorphing}
\tikzset{
  fn/.style   = {circle, draw=ele, fill=ele!12, line width=0.9pt,
                 minimum size=8mm, inner sep=0pt},
  fnwide/.style = {fn, rectangle, rounded corners=3.5mm, minimum width=17mm,
                   minimum height=8mm},
  coef/.style = {rectangle, draw=black!80, fill=black!6, line width=0.9pt,
                 minimum size=7mm, inner sep=0pt},
  coefwide/.style = {coef, minimum width=17mm, minimum height=7mm},
  fntall/.style = {fn, rectangle, rounded corners=3.5mm, minimum width=8mm,
                   minimum height=17mm},
  coeftall/.style = {coef, minimum width=7mm, minimum height=17mm},
  frame/.style = {draw=black!50, dashed, rounded corners=2mm, line width=0.7pt},
  cont/.style = {draw=ele, line width=0.9pt, decorate,
                 decoration={snake, amplitude=0.55mm, segment length=2.4mm,
                             pre length=0.4mm, post length=0.8mm}},
  disc/.style = {draw=black!80, line width=0.9pt},
  leg/.style  = {font=\small, inner sep=2pt},
}

%
% The convention: an expansion is a contraction, drawn as in MPS / SVD diagrams,
% except that a leg is one of two kinds.
%   cont  a wavy leg: a continuous argument, r or r_1 (a function of position)
%   disc  a plain line: a finite index, i or mu (ordinary linear algebra)
% Nodes likewise:
%   fn    a circle, in the electron colour: a function of position (psi, phi, chi)
%   coef  a square: a plain array of numbers (C, c)
%   coefwide a wide square box: an array with several legs on one side
%         (C_ij, H_ij,kl)
%   fntall, coeftall  the same two turned upright, for diagrams read left to
%         right (legs leave from the left and right sides)
%   frame a dashed box around a group of nodes that contracts to one array
%   fnwide a wide rounded box: a function of several positions (psi(r1, r2)),
%         its cont legs dropped in parallel from the bottom edge
% Joining two legs means summing over that index, so
%   psi(r1, r2) = sum_ij C_ij phi_i(r1) phi_j(r2)
% is  fn with two cont legs  =  fn -disc- coef -disc- fn,  each fn keeping its
% cont leg. Label a leg with node[leg] at its free end (or mid-way for a
% contracted one, e.g. "i").
%
% Minimal use:
%   \begin{tikzpicture}
%     \node[fn] (p) at (0,0) {$\varphi$};
%     \draw[cont] (p) -- ++(0,-1) node[leg, below] {$\mathbf r$};
%     \draw[disc] (p) -- ++(-1,0) node[leg, left] {$i$};
%   \end{tikzpicture}
\usepackage{tikz}
\usetikzlibrary{decorations.pathmorphing}
\tikzset{
  fn/.style   = {circle, draw=ele, fill=ele!12, line width=0.9pt,
                 minimum size=8mm, inner sep=0pt},
  fnwide/.style = {fn, rectangle, rounded corners=3.5mm, minimum width=17mm,
                   minimum height=8mm},
  coef/.style = {rectangle, draw=black!80, fill=black!6, line width=0.9pt,
                 minimum size=7mm, inner sep=0pt},
  coefwide/.style = {coef, minimum width=17mm, minimum height=7mm},
  fntall/.style = {fn, rectangle, rounded corners=3.5mm, minimum width=8mm,
                   minimum height=17mm},
  coeftall/.style = {coef, minimum width=7mm, minimum height=17mm},
  frame/.style = {draw=black!50, dashed, rounded corners=2mm, line width=0.7pt},
  cont/.style = {draw=ele, line width=0.9pt, decorate,
                 decoration={snake, amplitude=0.55mm, segment length=2.4mm,
                             pre length=0.4mm, post length=0.8mm}},
  disc/.style = {draw=black!80, line width=0.9pt},
  leg/.style  = {font=\small, inner sep=2pt},
}

%
% The convention: an expansion is a contraction, drawn as in MPS / SVD diagrams,
% except that a leg is one of two kinds.
%   cont  a wavy leg: a continuous argument, r or r_1 (a function of position)
%   disc  a plain line: a finite index, i or mu (ordinary linear algebra)
% Nodes likewise:
%   fn    a circle, in the electron colour: a function of position (psi, phi, chi)
%   coef  a square: a plain array of numbers (C, c)
%   coefwide a wide square box: an array with several legs on one side
%         (C_ij, H_ij,kl)
%   fntall, coeftall  the same two turned upright, for diagrams read left to
%         right (legs leave from the left and right sides)
%   frame a dashed box around a group of nodes that contracts to one array
%   fnwide a wide rounded box: a function of several positions (psi(r1, r2)),
%         its cont legs dropped in parallel from the bottom edge
% Joining two legs means summing over that index, so
%   psi(r1, r2) = sum_ij C_ij phi_i(r1) phi_j(r2)
% is  fn with two cont legs  =  fn -disc- coef -disc- fn,  each fn keeping its
% cont leg. Label a leg with node[leg] at its free end (or mid-way for a
% contracted one, e.g. "i").
%
% Minimal use:
%   \begin{tikzpicture}
%     \node[fn] (p) at (0,0) {$\varphi$};
%     \draw[cont] (p) -- ++(0,-1) node[leg, below] {$\mathbf r$};
%     \draw[disc] (p) -- ++(-1,0) node[leg, left] {$i$};
%   \end{tikzpicture}
\usepackage{tikz}
\usetikzlibrary{decorations.pathmorphing}
\tikzset{
  fn/.style   = {circle, draw=ele, fill=ele!12, line width=0.9pt,
                 minimum size=8mm, inner sep=0pt},
  fnwide/.style = {fn, rectangle, rounded corners=3.5mm, minimum width=17mm,
                   minimum height=8mm},
  coef/.style = {rectangle, draw=black!80, fill=black!6, line width=0.9pt,
                 minimum size=7mm, inner sep=0pt},
  coefwide/.style = {coef, minimum width=17mm, minimum height=7mm},
  fntall/.style = {fn, rectangle, rounded corners=3.5mm, minimum width=8mm,
                   minimum height=17mm},
  coeftall/.style = {coef, minimum width=7mm, minimum height=17mm},
  frame/.style = {draw=black!50, dashed, rounded corners=2mm, line width=0.7pt},
  cont/.style = {draw=ele, line width=0.9pt, decorate,
                 decoration={snake, amplitude=0.55mm, segment length=2.4mm,
                             pre length=0.4mm, post length=0.8mm}},
  disc/.style = {draw=black!80, line width=0.9pt},
  leg/.style  = {font=\small, inner sep=2pt},
}
